\documentclass[11pt]{article}
\usepackage[margin=1in]{geometry}
\usepackage{amsmath,amssymb,amsthm,mathtools}
\usepackage{bm}
\usepackage{booktabs}
\usepackage{enumitem}
\usepackage{xcolor}
\usepackage{hyperref}
\hypersetup{colorlinks=true, linkcolor=blue!50!black}

% --- minimal standalone macros (names match the EFGP paper's) --------------
\newcommand{\E}{\mathbb{E}}
\newcommand{\R}{\mathbb{R}}
\newcommand{\Eqx}{\E_{q(x_i)}}
\newcommand{\Eqxx}{\E_{q(x_i,x_{i+1})}}
\newcommand{\Cov}{\operatorname{Cov}}
\newcommand{\diag}{\operatorname{diag}}
\newcommand{\tr}{\operatorname{tr}}
\newcommand{\GP}{\mathcal{GP}}
\newcommand{\given}{\,|\,}
\newcommand{\code}[1]{\texttt{\small #1}}
\emergencystretch=3em
\hbadness=10000

\title{Updating the latent posterior $q(x)$ with the drift posterior $q(f)$
fixed}
\author{}
\date{}

\begin{document}
\maketitle
\vspace{-2.5em}

% ==========================================================================
\section*{What this note is, in one paragraph}
% ==========================================================================

\noindent
We are fitting a latent stochastic differential equation: a hidden
trajectory $x_0,\dots,x_T$ evolves under an unknown drift $f$, and we see
only noisy observations $y$. Inference is variational EM. We keep a
Gaussian approximation $q(x)$ to the hidden path and a Gaussian process
posterior $q(f)$ to the drift, and we improve them in alternation. This
note covers exactly one of those alternating updates: \textbf{holding the
drift posterior $q(f)$ fixed, improve the path posterior $q(x)$.} The
punchline is that this update is one \emph{exact} rewrite followed by one
\emph{single} approximation --- and that one approximation, applied once,
determines every quantity the update needs.

\bigskip
\noindent\fbox{\parbox{\dimexpr\textwidth-2\fboxsep-2\fboxrule}{%
\textbf{The whole update in three moves.}
\begin{description}[leftmargin=2.6em,itemsep=2pt,topsep=3pt]
\item[\textmd{\textbf{M1 (exact).}}] Rewrite the objective so that all of
its dependence on the drift $f$ funnels through just \emph{three numbers}
per time step --- averages of the drift and its slope over $q(x_i)$.
\item[\textmd{\textbf{M2 (bookkeeping).}}] We compute those three numbers
fast, in a way that hides them from automatic differentiation. So we must
tell the optimiser, by hand, how each of them changes with $q(x)$.
\item[\textmd{\textbf{M3 (the one approximation).}}] To get those
derivatives we make a single assumption --- \emph{the drift is a straight
line across the width of $q(x_i)$} --- and turn the crank. Every
derivative the update needs falls out of that one assumption.
\end{description}}}

% ==========================================================================
\section{Setup, and what a \texorpdfstring{$q(x)$}{q(x)} step must compute}
% ==========================================================================

\paragraph{The model.}
Time is discretised (Euler--Maruyama), so between samples the latent state
takes a noisy step driven by the drift:
\begin{equation}
x_{i+1} = x_i + \Delta_i\, f(x_i) + \epsilon_i,
\qquad \epsilon_i \sim \mathcal N(0,\Delta_i\Sigma).
\label{eq:euler}
\end{equation}
The variational posterior over the path is a Gaussian Markov chain: each
state has a marginal $q(x_i)=\mathcal N(m_i,S_i)$, and neighbours share a
cross-covariance $S_{i,i+1}=\Cov(x_i,x_{i+1})$. The drift posterior
$q(f)$ is already in hand from the other EM update; all we use from it is
the posterior-\emph{mean} drift, written $\bar f(x)$. (In the EFGP basis
$\bar f_r(x)=\phi(x)^*\mu_r$, but nothing below needs that form.)

\paragraph{Why we only ever compute one gradient.}
We improve $q(x)$ by natural-gradient ascent on the ELBO. For a Gaussian
posterior there is a useful shortcut: writing $q(x)$ in its
\emph{mean parameters}
\begin{equation}
\mu = \bigl(\E[x_i],\ \E[x_ix_i^\top],\ \E[x_ix_{i+1}^\top]\bigr),
\end{equation}
the natural gradient is just the \emph{ordinary} gradient in $\mu$,
\begin{equation}
\widetilde\nabla\,\mathcal L \;=\; \partial\mathcal L/\partial\mu .
\label{eq:duality}
\end{equation}
No Fisher matrix, no special calculus: we plainly differentiate the ELBO
with respect to $\mu$, damp the result, and step. That machinery --- the
damping, the conversion back to $(m_i,S_i,S_{i,i+1})$ --- is the SING
backbone and is not our concern here. \textbf{Our only job is to hand the
backbone $\partial\mathcal L/\partial\mu$.}

\paragraph{Only one piece of the ELBO involves the drift.}
The ELBO splits into four groups,
\begin{equation}
\mathcal L \;=\;
\underbrace{\E_{q(x)}[\log p(y\given x)]}_{\text{observations}}
\;+\; \sum_i \underbrace{\E_{q(x,f)}[\log p(x_{i+1}\given x_i,f)]}_{\text{transitions}}
\;-\; \underbrace{\mathrm{KL}[q(f)\,\|\,p(f)]}_{\text{no }q(x)}
\;+\; \underbrace{H[q(x)] + \E[\log p(x_0)]}_{\text{entropy, initial}} .
\end{equation}
The observation, entropy, and initial terms are ordinary Gaussian
expressions in $q(x)$; the backbone differentiates them directly. The KL
does not involve $q(x)$ at all. \textbf{Only the transition term couples
$q(x)$ to the drift}, so it is the only term we have to work on.
Everything from here concerns a single transition $i$; call its
contribution to the objective $\code{trm}_i$.

% ==========================================================================
\section{M1: reduce the transition term to three drift moments (exact)}
% ==========================================================================

Start from the log-transition of \eqref{eq:euler} and take its average
under the drift posterior $q(f)$. Two clean things happen. First, the
average of $f$ is the posterior mean $\bar f$. Second, the average of
$f^\top\Sigma^{-1}f$ is $\bar f^\top\Sigma^{-1}\bar f$ plus a variance
term $V$ (the drift's own uncertainty); we carry $V$ along and dispose of
it in M3. Writing $\delta_i := x_{i+1}-x_i$,
\begin{equation}
\E_{q(f)}\bigl[\log p(x_{i+1}\given x_i,f)\bigr]
= -\frac{\delta_i^\top\Sigma^{-1}\delta_i}{2\Delta_i}
+ \delta_i^\top\Sigma^{-1}\bar f(x_i)
- \frac{\Delta_i}{2}\Bigl[\bar f(x_i)^\top\Sigma^{-1}\bar f(x_i) + V(x_i)\Bigr]
+ \text{const}.
\label{eq:qf-averaged}
\end{equation}
Now average over the path, $\Eqxx[\cdot]$. Two exact simplifications
leave the whole thing depending on $q(x)$ through only a handful of
quantities.

\paragraph{The first term is already elementary.}
$\Eqxx[\delta_i^\top\Sigma^{-1}\delta_i]$ is a Gaussian average of a
quadratic, so it expands into traces of $\E[x_ix_i^\top]$,
$\E[x_{i+1}x_{i+1}^\top]$, $\E[x_ix_{i+1}^\top]$ and the means. It is
\emph{linear in $\mu$}: no drift, no approximation, and its
$\mu$-gradient is immediate.

\paragraph{The middle term is collapsed by Stein's lemma.}
The cross term $\Eqxx[\delta_i^\top\Sigma^{-1}\bar f(x_i)]$ couples $x_i$
and $x_{i+1}$. Stein's lemma trades that coupling for two \emph{marginal}
averages against known $q(x)$ statistics:
\begin{equation}
\Eqxx\bigl[\delta_i^\top\Sigma^{-1}\bar f(x_i)\bigr]
= d_i^\top\Sigma^{-1}\,\Eqx[\bar f]
\;+\; \tr\!\bigl[\Sigma^{-1}\,\Eqx[J_{\bar f}]\,C_i^\top\bigr],
\qquad
\begin{aligned}
d_i &:= m_{i+1}-m_i,\\
C_i &:= S_{i,i+1}-S_i,
\end{aligned}
\label{eq:stein}
\end{equation}
where $J_{\bar f}$ is the Jacobian (the slope) of the mean drift.

\paragraph{What remains: three drift moments.}
After \eqref{eq:qf-averaged}--\eqref{eq:stein}, the transition term is a
plain algebraic formula in $(m_i,S_i,S_{i,i+1})$ and just three averages
of the drift over the marginal $q(x_i)$:
\begin{equation}
\boxed{\;
\begin{aligned}
\mathrm{Ef}_i &:= \Eqx[\bar f] && (\text{mean drift}), &
\mathrm{Edfdx}_i &:= \Eqx[J_{\bar f}] && (\text{mean slope}), \\
\mathrm{Eff}_i &:= \Eqx\bigl[\bar f^\top\Sigma^{-1}\bar f\bigr] && (\text{mean square}).
\end{aligned}
\;}
\label{eq:moments}
\end{equation}
Concretely, dropping the constant and the overall $-1/(2\Delta_i\sigma^2)$
prefactor, one transition contributes
\begin{equation}
\code{trm}_i \;\propto\;
\Eqxx\bigl[\delta_i^\top\Sigma^{-1}\delta_i\bigr]
\;+\;\Delta_i^2\,\mathrm{Eff}_i
\;-\;2\Delta_i\Bigl(d_i^\top\Sigma^{-1}\,\mathrm{Ef}_i
   + \tr\bigl[\Sigma^{-1}\,\mathrm{Edfdx}_i\,C_i^\top\bigr]\Bigr).
\label{eq:assembled}
\end{equation}

\medskip\noindent
\textbf{This is M1, and it is exact.} It is also the entire interface
between the drift model and the path update: give the backbone the three
numbers in \eqref{eq:moments}, and it assembles \eqref{eq:assembled}
knowing nothing about where they came from. This mirrors the other EM
update, where uncertainty in the path enters the drift fit only by
replacing fixed sufficient statistics with their $q(x)$-averages; here it
is the same idea, read the other way.

% ==========================================================================
\section{M2: the moments' derivatives must be supplied by hand}
% ==========================================================================

We need $\partial\,\code{trm}_i/\partial\mu$. The explicit algebra of
\eqref{eq:assembled} --- the $\delta$-quadratic and the $d_i,C_i$ pieces
--- differentiates by ordinary autodiff, exactly. The difficulty is that
the three moments \eqref{eq:moments} are \emph{themselves} functions of
$(m_i,S_i)$, because they are averages over $q(x_i)$. So the chain rule
also demands their derivatives:
\begin{equation}
\frac{\partial\,\mathrm{Ef}_i}{\partial(m_i,S_i)},\qquad
\frac{\partial\,\mathrm{Edfdx}_i}{\partial(m_i,S_i)},\qquad
\frac{\partial\,\mathrm{Eff}_i}{\partial(m_i,S_i)}
\qquad\text{(six derivatives).}
\label{eq:six-slots}
\end{equation}

\paragraph{Why autodiff cannot just do it.}
Each moment has a closed form (a Gaussian average of a complex
exponential), and for speed we evaluate all $T$ of them together in
\emph{one} batched transform. But a derivative is needed
\emph{per transition}. Asking autodiff to differentiate that batched
transform in place would pay one full transform \emph{per transition}
instead of one for all of them --- orders of magnitude slower. So we
compute the moments off the autodiff graph, pass them in as constants,
and attach the six derivatives by hand (a \code{custom\_vjp} in JAX).
Autodiff then differentiates \eqref{eq:assembled} normally and picks up
our injected derivatives through the chain rule.

\paragraph{What the exact derivatives are.}
For a Gaussian $q(x_i)=\mathcal N(m_i,S_i)$, the derivatives of an average
are given by Price's theorem (its mean case is Bonnet's lemma):
\begin{equation}
\partial_{m_i}\Eqx[g] = \Eqx[\nabla g],
\qquad
\partial_{S_i}\Eqx[g] = \tfrac12\,\Eqx[\nabla^2 g].
\label{eq:price}
\end{equation}
That is: differentiating in the mean pulls down a \emph{gradient} of the
integrand; differentiating in the covariance pulls down a \emph{Hessian}.
Applied to the three moments, the six exact answers are
\begin{center}
\renewcommand{\arraystretch}{1.5}\small
\begin{tabular}{@{}p{3.0cm} p{4.2cm} p{7.4cm}@{}}
\toprule
moment & exact $\partial/\partial m_i$ & exact $\partial/\partial S_i$ \\
\midrule
$\Eqx[\bar f]$
  & $\Eqx[J_{\bar f}]$
  & $\tfrac12\,\Eqx[H_{\bar f}]$ \\
$\Eqx[J_{\bar f}]$
  & $\Eqx[H_{\bar f}]$
  & $\tfrac12\,\Eqx[\nabla^2 J_{\bar f}]$ \\
$\Eqx[\bar f^\top\Sigma^{-1}\bar f]$
  & $2\,\Eqx[J_{\bar f}^\top\Sigma^{-1}\bar f]$
  & $\Eqx[J_{\bar f}^\top\Sigma^{-1}J_{\bar f}]
    + \sum_r\sigma_r^{-2}\,\Eqx[\bar f_r H_{\bar f_r}]$ \\
\bottomrule
\end{tabular}
\end{center}
where $H_{\bar f}$ is the Hessian (the curvature) of the mean drift. Five
of the six answers need this curvature; only the top-left --- the slope
$\Eqx[J_{\bar f}]$, which we already have --- does not. Computing the
curvature everywhere is the expense we would like to avoid. M3 avoids it.

% ==========================================================================
\section{M3: one approximation, and the whole table falls out}
% ==========================================================================

Here is the single assumption behind the entire $q(x)$ update:

\begin{quote}\itshape
Across the width of $q(x_i)$, replace the drift by a straight line.
\end{quote}

\noindent
We fix the line by \emph{statistical linearisation}: the affine map that
matches $\bar f$'s first two moments under $q(x_i)$,
\begin{equation}
\ell(x) \;=\; a + B\,(x-m_i),
\qquad
a := \Eqx[\bar f],
\qquad
B := \Eqx[J_{\bar f}],
\label{eq:linearise}
\end{equation}
its offset the mean drift and its slope the \emph{expected} Jacobian. This
is not the Taylor tangent at $m_i$: by Stein's lemma --- the identity
already used in \eqref{eq:stein} --- $\Eqx[J_{\bar f}]$ equals the
least-squares slope $\Cov(\bar f,x)\,S_i^{-1}$, so $\ell$ is the best
affine fit to $\bar f$ over the whole blob (statistical linearisation; S\"arkk\"a, 2013; Opper \& Archambeau, 2009). Both coefficients $a,B$ are
the exact gmix moments of M1; the slope is $\Eqx[J_{\bar f}]$
throughout, \emph{never} a point value $J_{\bar f}(m_i)$. We use $\ell$ in
place of $\bar f$ inside every $q(x_i)$-average; the six derivatives then
come out with no further choices.

\paragraph{How to picture it.}
The assumption is not that the drift is simple --- we linearise afresh at
\emph{every} $m_i$, so the tangents track the true curved field along the
path. It is a statement about $q(x)$: each marginal $q(x_i)$ is narrow
enough that the drift is effectively straight across that one blob. Take a
one-dimensional bistable drift, $f(x)=x-x^3$. Over a Gaussian $q(x_i)=\mathcal N(m_i,S_i)$ of width $\sqrt{S_i}$ we
pretend $f$ is a straight line (of slope $\Eqx[f']$; the point tangent
$f'(m_i)=1-3m_i^2$ differs from it only at the order below). The mismatch
is set by the curvature, $f''=-6x$: the error in any averaged drift
quantity is of order $\tfrac12 f''\,S_i$, i.e.\ (curvature) $\times$
(width).
It is zero if $f$ is genuinely affine (a linear SDE, a line attractor),
and small whenever $q(x_i)$ is narrow \emph{relative to the scale on which
the drift bends} --- the length-scale $\ell$ --- so the honest condition
is $\|S_i\|\ll\ell^2$, not $\|S_i\|$ small in absolute terms. We expand at
the mean precisely because there the first-order error cancels, leaving
only this second-order term.

\paragraph{When it breaks.}
The error is largest where a \emph{wide} blob sits on a \emph{sharply
bending} part of the field --- both factors at once. In practice that
means: the cold-start iterations, before the smoother has localised
$q(x)$; stretches of the trajectory with few or noisy observations, where
$q(x_i)$ stays broad; and sharp turns of the flow (a fast reversal, a
separatrix) where $\|H_{\bar f}\|$ is large. A tight blob on a gentle part
of the field, or any blob on a near-linear part, is fine.

\paragraph{Why the breaking case does not hurt much.}
A wide $q(x_i)$ is, by definition, a latent state the data do not pin
down. A slightly mis-directed gradient on a weakly-determined coordinate
has little leverage on the reconstruction, and SING's $\rho$-damping
attenuates any single overconfident step on top of that --- so the
approximation is worst exactly where the answer is least sensitive to it.
It is also self-correcting: in well-observed regions the update tightens
$q(x_i)$, which makes the next linearisation better, so error shrinks over
EM iterations rather than accumulating. And the damage is bounded even
in the transient: the moment \emph{values} fed to the backbone are the
true Gaussian averages of \eqref{eq:moments} (only their \emph{derivatives}
assume a straight line), so the step's base point stays exact and only its
\emph{direction} is affected; and each transition's error is local to its
own term in the sum, not compounded along the path. (These are the same
regimes flagged below for optionally restoring the $V$ term.)

\paragraph{And the worst case is not one we could trust anyway.}
Sharpen the failure condition: the error is order-one only once the
positional standard deviation approaches the length-scale,
$\sqrt{\|S_i\|}\gtrsim\ell$. But $\ell$ \emph{is} the scale on which the
drift varies, and the only information we have about the drift at that
scale comes from watching $x$ move between states we have localised to
precision $\sqrt{\|S_i\|}$. If $\sqrt{\|S_i\|}\gg\ell$, each look at the
drift is a smear wider than the feature it would resolve, so the data
cannot have identified structure at scale $\ell$ in the first place: a
confident $\hat\ell$ far below the positional spread is simply
unsupported. Where the linearisation is worst, the length-scale it
conditions on is itself untrustworthy --- the two degrade together, and
neither is a regime one would report a fit from. Conversely, any
well-identified converged fit has $\sqrt{\|S_i\|}\ll\ell$, the path pinned
to well within the drift's variation scale, which is exactly where
linearisation is excellent; validity and convergence coincide. (This is
the typical, whole-trajectory statement --- a few wide-$S_i$ transitions
in an otherwise sharp fit are just the weakly-identified states above. And
since only $\tr(H_{\bar f}S_i)$ enters, uncertainty along directions in
which the drift is flat costs nothing, so the true condition is even more
forgiving than $\|S_i\|<\ell^2$.)

\paragraph{The six gradients.}
$\ell$ is affine, so its moments under $q(x_i)=\mathcal N(m,S)$ are exact
(the linear part averages out; the Jacobian is the constant $B$; the
quadratic adds one trace),
\[
\Eqx[\ell]=a,\qquad
\Eqx[J_\ell]=B,\qquad
\Eqx[\ell^\top\Sigma^{-1}\ell]=a^\top\Sigma^{-1}a+\tr(B^\top\Sigma^{-1}B\,S),
\]
with the $m$-dependence carried only by $\Eqx[\ell]=a+B(m-m_i)$.
Differentiating in $(m,S)$ at $m=m_i$ gives the six injected gradients:
\begin{align*}
\partial_m\Eqx[\bar f] &= B=\Eqx[J_{\bar f}], &
\partial_S\Eqx[\bar f] &= 0,\\
\partial_m\Eqx[J_{\bar f}] &= 0, &
\partial_S\Eqx[J_{\bar f}] &= 0,\\
\partial_m\Eqx[\bar f^\top\Sigma^{-1}\bar f] &= 2B^\top\Sigma^{-1}a, &
\partial_S\Eqx[\bar f^\top\Sigma^{-1}\bar f] &= B^\top\Sigma^{-1}B \;\succeq 0,
\end{align*}
$a=\Eqx[\bar f]$, $B=\Eqx[J_{\bar f}]$ \eqref{eq:linearise}, both exact
gmix moments. The top-left is the exact Price gradient
$\partial_m\Eqx[\bar f]=\Eqx[J_{\bar f}]$; the bottom-right
$B^\top\Sigma^{-1}B\succeq0$ is the Gauss--Newton curvature of the sum of
squares $\bar f^\top\Sigma^{-1}\bar f$ (Khan \& Lin, 2017), keeping the
precision update \eqref{eq:duality} PSD.

\paragraph{The drift variance $V$: one more term of the same kind, dropped by default.}
$V_r(x)=\sigma_r^{-2}\phi(x)^*A_r^{-1}\phi(x)$ is a quadratic form in the
Fourier features --- like the mean-square $\bar f_r^2=\phi^*\mu_r\mu_r^*\phi$
--- so it too is a sum of difference-frequency modes
$\sum_\delta\omega_r(\delta)e^{2\pi i\delta^\top x}$, and its
$q(x_i)$-average is the \emph{same} characteristic-function envelope-NUFFT
as every other term (value and $m$-gradient exact, $S$-slot dropped). The
dense $A_r^{-1}$ only sets the coefficients $\omega_r(\delta)$: (i) compute
the drift posterior variance by any method --- Hutchinson on $A_r^{-1}$, or
the pathwise/Matheron estimator we use for UQ --- then (ii) apply the same
gmix envelope. Production omits $V$ altogether: $\nabla V\approx0$ wherever
the data surround $x$ evenly at scale $\ell$
($\|\nabla\Eqx[V]\|\approx5\times10^{-5}\|\partial_m\mathrm{Eff}\|$ on our
benchmark, recovery unchanged to 3--4 digits). Restore it (opt-in) when
that flatness fails --- clustered data, sharp boundary turns, or cold
start.

% ==========================================================================
\section{One inner iteration, end to end}
% ==========================================================================

\begin{enumerate}[leftmargin=2.2em,itemsep=2pt]
\item From the current $q(x)$, read off $m_i,S_i,S_{i,i+1}$ and form
  $d_i=m_{i+1}-m_i$, $C_i=S_{i,i+1}-S_i$.
\item Update the drift posterior $q(f)$ (the other EM step), giving the
  mean drift $\bar f$.
\item \textbf{[M1]} Evaluate the three moments \eqref{eq:moments} at all
  transitions in one batched transform.
\item \textbf{[M2, M3]} Hand them to the backbone as constants, with the
  linearised derivative table of \S4 attached.
\item Autodiff the assembled scalar \eqref{eq:assembled} in $\mu$; by
  \eqref{eq:duality} this is the natural gradient. Damp and step.
\item Repeat for $n_{\mathrm{estep}}$ inner iterations, then take the
  kernel (hyperparameter) M-step.
\end{enumerate}

% ==========================================================================
\section*{Crosswalk to the long note}
% ==========================================================================

\noindent
\begin{tabular}{@{}l l@{}}
\toprule
here & \texttt{efgp\_estep.tex} \\
\midrule
\S1 setup, what the step computes & \S2 (interface), \S6 setup \\
\S2 M1, three drift moments (exact) & \S6 ``the transition term'', items (i)--(ii) \\
\S3 M2, hand-supplied derivatives & \S6 ``why we can't just autodiff'', App.~C \\
\S4 M3, the one approximation & \S6 Gauss--Newton framework, items (iii)--(iv) \\
\S4 the $V$ drop & \S6 item (iv); App.~B (restoration) \\
how the moments are computed & \S4 (three backends), \texttt{efgp\_estep\_charfun.tex} \\
\bottomrule
\end{tabular}

\end{document}
